\section*{Capítulo \ref{ch:g9:algebra} --- Álgebra e equações}

\begin{solution}{exo:g9:algebra:1}
$4(3x - 2) = 12x - 8$.

$(x+6)(x+2) = x^2 + 2x + 6x + 12 = x^2 + 8x + 12$.

$(2x-3)(x+5) = 2x^2 + 10x - 3x - 15 = 2x^2 + 7x - 15$.

$(x-1)(x-9) = x^2 - 9x - x + 9 = x^2 - 10x + 9$.
\end{solution}

\begin{solution}{exo:g9:algebra:2}
$(x+8)^2 = x^2 + 16x + 64$.

$(3x-4)^2 = 9x^2 - 24x + 16$.

$(x+10)(x-10) = x^2 - 100$.

$(5x+2)(5x-2) = 25x^2 - 4$.
\end{solution}

\begin{solution}{exo:g9:algebra:3}
$101^2 = (100+1)^2 = 10000 + 200 + 1 = 10201$.

$59^2 = (60-1)^2 = 3600 - 120 + 1 = 3481$.

$53 \times 47 = (50+3)(50-3) = 2500 - 9 = 2491$.
\end{solution}

\begin{solution}{exo:g9:algebra:4}
$9x + 12 = 3(3x + 4)$.

$x^2 - 64 = (x+8)(x-8)$.

$x^2 + 14x + 49 = (x+7)^2$.

$18x^2 - 6x = 6x(3x - 1)$.
\end{solution}

\begin{solution}{exo:g9:algebra:5}
$4x + 7 = 19$: $4x = 12$, então $x = 3$.

$6x - 5 = 2x + 11$: $4x = 16$, então $x = 4$.

$5(x-3) = 3x + 1$: $5x - 15 = 3x + 1$, então $2x = 16$ e $x = 8$.
\end{solution}

\begin{solution}{exo:g9:algebra:6}
$(x+4)(3x-12) = 0$: $x = -4$ ou $x = 4$.

$x^2 - 121 = (x-11)(x+11) = 0$: $x = 11$ ou $x = -11$.

$(2x-1)^2 = 0$: a solução única $x = \frac12$.

$x^2 = 8x$: $x(x - 8) = 0$, então $x = 0$ ou $x = 8$.
\end{solution}

\begin{solution}{exo:g9:algebra:7}
\omterm{ex:g1:subtraction:difference}{Diferença} de quadrados com $a = 3x + 2$ e $b = 5$:
\[
E = (3x + 2)^2 - 5^2 = (3x + 2 + 5)(3x + 2 - 5) = (3x + 7)(3x - 3).
\]
Regra do \omterm{def:g2:mult:def}{produto} nulo: $3x + 7 = 0$ dá $x = -\frac73$, e
$3x - 3 = 0$ dá $x = 1$.
\end{solution}

\begin{solution}{exo:g9:algebra:8}
$3x - 4 \geq 8$: $3x \geq 12$, então $x \geq 4$ (bolinha cheia em $4$,
traço para a direita).

$5 - 2x > 1$: $-2x > -4$, e dividir por $-2$ inverte o sinal: $x < 2$
(bolinha vazada em $2$, traço para a esquerda).

$4(x+1) \leq x - 5$: $4x + 4 \leq x - 5$, então $3x \leq -9$, logo
$x \leq -3$ (bolinha cheia em $-3$, traço para a esquerda).
\end{solution}

\begin{solution}{exo:g9:algebra:9}
Seja $x$ o número: $3x + 7 = 5x - 9$. Então $16 = 2x$, logo $x = 8$.
Conferindo: $3 \times 8 + 7 = 31$ e $5 \times 8 - 9 = 31$.
\end{solution}

\begin{solution}{exo:g9:algebra:10}
\emph{1.} $(n+1)(n-1) = n^2 - 1$ (terceiro \omterm{def:g2:mult:def}{produto} notável). Os
vizinhos de $n$ são $n - 1$ e $n + 1$, então o \omterm{def:g2:mult:def}{produto} deles é
$n^2 - 1$: por exemplo, $24 \times 26 = 625 - 1 = 624$.

\emph{2.} Três inteiros consecutivos podem ser escritos $n - 1$, $n$,
$n + 1$. A \omterm{def:g1:addition:def}{soma} deles é
\[
(n-1) + n + (n+1) = 3n ,
\]
\omterm{def:g9:arith:divisor}{múltiplo} de $3$ (três vezes o do meio).
\end{solution}

\begin{solution}{pb:g9:algebra:1}
\textbf{1.} $(10a + b) - (10b + a) = 9a - 9b = 9(a - b)$: a \omterm{ex:g1:subtraction:difference}{diferença}
entre um número de dois algarismos e o invertido dele é sempre múltipla
de $9$. Conferindo: $72 - 27 = 45 = 9 \times 5$, e
$a - b = 7 - 2 = 5$.

\textbf{2.} Desenvolvendo o lado direito:
$9(11a + b) + a + b + c = 99a + 9b + a + b + c = 100a + 10b + c$.
Então o número é igual a um \omterm{def:g9:arith:divisor}{múltiplo} de $9$ mais a \omterm{def:g1:addition:def}{soma} dos algarismos
dele: dividindo por $9$, os dois deixam o mesmo \omterm{def:g3:division:remainder}{resto}. Em particular,
\omterm{def:g3:division:remainder}{resto} $0$ para um significa \omterm{def:g3:division:remainder}{resto} $0$ para o outro: o critério do nove,
demonstrado.

\textbf{3.} $9(11a + b)$ também é \omterm{def:g9:arith:divisor}{múltiplo} de $3$, então a mesma
identidade mostra que o número e a \omterm{def:g1:addition:def}{soma} dos algarismos deixam o mesmo
\omterm{def:g3:division:remainder}{resto} na divisão por $3$: o critério do três.

\textbf{4.} Escreva $47 = 9k + 2$ e $86 = 9l + 5$. Então
\[
47 \times 86 = (9k + 2)(9l + 5)
= 9\,(9kl + 5k + 2l) + 10 :
\]
o \omterm{def:g2:mult:def}{produto} deixa o mesmo \omterm{def:g3:division:remainder}{resto} que $2 \times 5 = 10$, ou seja, $1$. Um
resultado correto tem, portanto, de deixar também \omterm{def:g3:division:remainder}{resto} $1$ --- o que
$4\,042$ deixa. Mas o teste só enxerga \omterm{def:g3:division:remainder}{restos}: $4\,042$ e $4\,402$
(algarismos trocados de lugar) passam igualzinho, então passar não
demonstra nada --- falhar é que condena.

\textbf{5.} Desenvolvendo:
$11(91a + 9b + c) - a + b - c + d = 1001a + 99b + 11c - a + b - c + d
= 1000a + 100b + 10c + d$. Então o número é \omterm{def:g6:wholes:divisible}{divisível} por $11$
exatamente quando a \omterm{def:g1:addition:def}{soma} alternada $-a + b - c + d$ é. Para $2\,926$:
$-2 + 9 - 2 + 6 = 11$, \omterm{def:g6:wholes:divisible}{divisível} por $11$: sim,
$2\,926 = 11 \times 266$.

\textbf{6.} $742 - 247 = 495$; $495 + 594 = 1\,089$. Qualquer exemplo
válido também cai em $1\,089$.

\textbf{7.} $(100a + 10b + c) - (100c + 10b + a) = 99(a - c)$. Para
$a - c = 2, 3, \dots, 9$:
\[
198,\ 297,\ 396,\ 495,\ 594,\ 693,\ 792,\ 891 .
\]

\textbf{8.} $100(m - 1) + 90 + (10 - m) = 100m - 100 + 90 + 10 - m
= 99m$. Então os três algarismos de $99m$ são $m - 1$, depois $9$,
depois $10 - m$ (confira em $495$: $m = 5$ dá $4$, $9$, $5$).

\textbf{9.} O invertido de $99m$ tem algarismos $10 - m$, $9$,
$m - 1$, então vale $100(10 - m) + 90 + (m - 1)$. Somando:
\[
100\,(m - 1 + 10 - m) + 180 + (10 - m + m - 1)
= 100 \times 9 + 180 + 9 = 1\,089 ,
\]
independentemente de $m$. Seja qual for o número de partida, o mágico
está seguro.

\textbf{10.} Se $a = c$, a subtração dá $0$ e o truque morre na hora.
Se $a - c = 1$, ela dá $99$, que precisa ser tratado como a cadeia de
três algarismos $099$ (invertida, $990$; $99 + 990 = 1\,089$: salvo!).
As letras miúdas: \omterm{ex:g1:subtraction:difference}{diferenças} de $1$ só funcionam se o espectador
escrever o resultado da subtração com três algarismos, \omterm{def:g1:counting:zero}{zero} à esquerda
incluído --- exigir uma \omterm{ex:g1:subtraction:difference}{diferença} de pelo menos $2$ evita a discussão.

\textbf{11.} $2(5m + 7) + d - 14 = 10m + 14 + d - 14 = 10m + d$: as
dezenas (e as centenas, de outubro a dezembro) mostram o mês; as
unidades (e as dezenas), o dia.

\textbf{12.} $2(5m + 9) + d = 10m + 18 + d$: subtraia $18$.

\textbf{13.} Pela questão 2,
$n - (\text{soma dos algarismos de } n) = 9(11a + b)$ para um $n$ de
três algarismos: sempre \omterm{def:g9:arith:divisor}{múltiplo} de $9$ (identidades parecidas
funcionam para qualquer quantidade de algarismos). Um \omterm{def:g9:arith:divisor}{múltiplo} de $9$
tem \omterm{def:g1:addition:def}{soma} dos algarismos múltipla de $9$ (questão 2 de novo), então os
algarismos que o espectador lê têm de ser completados até o \omterm{def:g9:arith:divisor}{múltiplo} de
$9$ seguinte: o algarismo que falta é esse complemento. Ambiguidade: se
os algarismos lidos já somam um \omterm{def:g9:arith:divisor}{múltiplo} de $9$, o algarismo riscado
poderia ser $0$ ou $9$ --- excluir o $0$ tira a dúvida.

\textbf{14.} $(10t + 5)^2 = 100t^2 + 100t + 25 = 100\,t(t + 1) + 25$:
escreva $t(t+1)$ e cole o $25$ atrás. $45^2$: $4 \times 5 = 20$, então
$2\,025$. $85^2$: $8 \times 9 = 72$, então $7\,225$. $105^2$:
$10 \times 11 = 110$, então $11\,025$.

\textbf{15.} Cinco consecutivos:
$n + (n+1) + (n+2) + (n+3) + (n+4) = 5n + 10 = 5(n + 2)$: \omterm{def:g9:arith:divisor}{múltiplo} de
$5$, sempre. Quatro consecutivos:
$n + (n+1) + (n+2) + (n+3) = 4n + 6 = 4(n + 1) + 2$: \omterm{def:g3:division:remainder}{resto} $2$ na
divisão por $4$, nunca \omterm{def:g9:arith:divisor}{múltiplo}. A álgebra converte ``parece
funcionar'' em \emph{sempre}, e ``não achei exemplo nenhum'' em
\emph{nunca} --- duas palavras que quantidade nenhuma de testes alcança
(\cref{met:g7:literal:test}).
\end{solution}
